\documentclass[10pt]{article}
\usepackage[T1]{fontenc}
\usepackage[utf8]{inputenc}
\usepackage[french]{babel}
\usepackage{amsmath,amssymb,mathrsfs,mathpazo}
\usepackage{geometry}
\geometry{a4paper,margin=2cm}
\usepackage{graphicx,array,tabularx,tabularray,colortbl,multirow,multicol,booktabs,enumitem,pifont,ulem,fancybox,diagbox,cancel}
\usepackage{pstricks,pstricks-add,pst-plot,pst-tree,pst-node,pst-text,pst-func,pst-math,pst-bspline,pst-3dplot,pst-eucl,pst-solides3d}
\let\pscalculate\undefined
\usepackage{pst-bezier}
\usepackage{tikz,pgfplots,tkz-tab}
\usetikzlibrary{arrows,arrows.meta,patterns,calc,positioning,trees,decorations.markings,shapes.geometric,shapes.misc,shapes.arrows,intersections,decorations.pathmorphing,decorations.pathreplacing,backgrounds,shadows,babel}
\pgfplotsset{compat=1.18}
\usepgfplotslibrary{fillbetween}
\usepackage[np]{numprint}
\usepackage{listings,makecell,pas-tableur,twoopt}
\let\barre\undefined
\usepackage[pstricks]{tablvar}
\definecolor{darkgreen}{RGB}{0,100,0}
\providecommand{\red}{\color{red}}
\providecommand{\blue}{\color{blue}}
\providecommand{\green}{\color{green}}
\providecommand{\black}{\color{black}}
\providecommand{\white}{\color{white}}
\providecommand{\gray}{\color{gray}}
\providecommand{\cyan}{\color{cyan}}
\providecommand{\magenta}{\color{magenta}}
\providecommand{\yellow}{\color{yellow}}
\providecommand{\violet}{\color{violet}}
\providecommand{\orange}{\color{orange}}
\providecommand{\euro}{\texteuro}
\providecommand{\EUR}{\texteuro}
\providecommand{\ieme}{\textsuperscript{e}}
\providecommand{\degres}{\ensuremath{^\circ}}
\providecommand{\degree}{\ensuremath{^\circ}}
\providecommand{\vv}[1]{\overrightarrow{#1}}
\providecommand{\Coord}[1]{#1}
\providecommand{\decosix}{}
\providecommand{\starredbullet}{\ensuremath{\bullet}}
\pagestyle{empty}
\setlength{\parindent}{0pt}
\usepackage{adjustbox}\DeclareUnicodeCharacter{2236}{:}\DeclareUnicodeCharacter{2190}{\ensuremath{\leftarrow}}
\providecommand{\R}{\mathbb{R}}
\providecommand{\N}{\mathbb{N}}
\providecommand{\Z}{\mathbb{Z}}
\providecommand{\Q}{\mathbb{Q}}
\providecommand{\D}{\mathbb{D}}
\providecommand{\C}{\mathbb{C}}
\providecommand{\e}{{\mathrm{e}}}
\providecommand{\vect}[1]{\overrightarrow{#1}}
\providecommand{\vectt}[1]{\overrightarrow{\mathrm{#1}}}
\providecommand{\vvt}[1]{\overrightarrow{\mathrm{#1}}}
\providecommand{\np}[2][]{{#2~#1}}
\providecommand{\vv}[1]{{\overrightarrow{#1}}}
\renewcommand{\d}{{\mathrm{d}}}
\providecommand{\degres}{{^\circ}}
\providecommand{\degre}{{^\circ}}
\providecommand{\blue}{\color{blue}}
\providecommand{\red}{\color{red}}
\providecommand{\enskip}{\quad}
\providecommand{\llbracket}{\lbrack\!\lbrack}
\providecommand{\rrbracket}{\rbrack\!\rbrack}
\providecommand{\euro}{{\text{€}}}
\providecommand{\no}{{\text{n°}}}
\providecommand{\interval}[2]{{[#1;#2]}}
\providecommand{\intervFO}[2]{{[#1;#2[}}
\providecommand{\dotfill}{\ldots}

\providecommand{\ds}{\displaystyle}
\providecommand{\textding}[1]{#1}
\providecommand{\Oij}{\(\mathrm{O};\vec{\imath},\vec{\jmath}\)}
\providecommand{\Oijk}{\(\mathrm{O};\vec{\imath},\vec{\jmath},\vec{k}\)}
\providecommand{\Ouv}{\(\mathrm{O};\vec{u},\vec{v}\)}
\providecommand{\intervFF}[2]{\ensuremath{[#1;#2]}}
\renewcommand{\euro}{\ifmmode\text{\texteuro}\else\texteuro\fi}
\DeclareUnicodeCharacter{00B0}{\ifmmode^{\circ}\else\textdegree\fi}
\DeclareUnicodeCharacter{20AC}{\ifmmode\text{\texteuro}\else\texteuro\fi}
% Énoncés : Éducation nationale. Transcription et corrigés : APMEP, Denis Vergès et contributeurs. Découpage et classement : CoopMaths, Sébastien Lozano, Éric Elter et contributeurs. Adaptation de la présentation : Maths Explorer.
% CC BY-SA : https://coopmaths.fr/www/about/
% https://coopmaths.fr/alea/static/e3c/2020/tex/e3c_2020_00_sujet28_3_cor.tex
\begin{document}
\noindent\textit{APMEP / CoopMaths. CC BY-SA. Présentation adaptée par Maths Explorer.}\par\medskip
\subsection*{1.}
	Diminuer de \(5\%\), c'est multiplier la quantité de déchets par \(1 - \dfrac{5}{100} = 1 - 0,05 = 0,95\). Donc l'entreprise s'engage à produire seulement en 2020 :
	\[
	6000 \times 0,95 = 5700 \text{ (tonnes)}.
	\]
	
	\subsection*{2.}
	
	\paragraph{a.} On a donc pour tout entier naturel \(n\),
	\[
	d_{n+1} = 0,95 d_n.
	\]
	
	\paragraph{b.} La relation précédente montre que la suite \((d_n)\) est une suite géométrique de raison \(0,95\) et de premier terme \(d_0 = 6000\).
	
	\paragraph{c.} On sait que pour tout entier naturel \(n\) :
	\[
	d_n = d_0 \times 0,95^n = 6000 \times 0,95^n.
	\]
	La quantité totale de déchets produite de 2019 à 2023 est :
	\begin{align*}
		T &= d_0 + d_1 + d_2 + d_3 + d_4 \\
		&= 6000 \times 0,95^0 + 6000 \times 0,95^1 + 6000 \times 0,95^2 + 6000 \times 0,95^3 + 6000 \times 0,95^4 \quad (1),
	\end{align*}
	d'où, par produit par \(0,95\) :
	\begin{align*}
		0,95T &= 6000 \times 0,95^1 + 6000 \times 0,95^2 + 6000 \times 0,95^3 + 6000 \times 0,95^4 + 6000 \times 0,95^5 \quad (2),
	\end{align*}
	d'où, par différence \((1) - (2)\) :
	\[
	0,05T = 6000 - 6000 \times 0,95^5,
	\]
	d'où on tire :
	\[
	T = \dfrac{6000 - 6000 \times 0,95^5}{0,05} \approx 27146 \text{ tonnes}.
	\]
	
	\subsection*{3.}
	\begin{tcolorbox}[colback=white,colframe=black,left=1mm,right=1mm,top=1mm,bottom=1mm,boxsep=0mm]
		\[
		\begin{array}{l}
			D \leftarrow 6000 \\
			N \leftarrow 0 \\
			\text{Tant que } D > 3600 \\
			\quad D \leftarrow 0,95 * D \\
			\quad N \leftarrow N + 1 \\
			\text{Fin Tant que}
		\end{array}
		\]
	\end{tcolorbox}
	
	L'algorithme donnera \(N = 10\), soit en 2029.

\end{document}
